\section{Checklist de Examen}

\begin{longtable}{p{0.24\linewidth}p{0.36\linewidth}p{0.31\linewidth}}
\toprule
Tema & Senal en el enunciado & Algoritmo que debes activar \\
\midrule
\endhead
Muestreo de aceptacion & Planes $(n,c)$, AQL, LPTD, $\alpha$, $\beta$, porcentaje defectuoso & Modelar $X\sim Bin(n,p)$ y calcular $P(X\le c)$. \\
Grafico $\bar X$ & Promedios muestrales, $\sigma$ conocida, nivel de confianza & Usar $\mu_0\pm z\sigma/\sqrt n$. \\
Grafico $\bar X/R$ & Subgrupos, rangos, constantes $A_2,D_3,D_4$ & Calcular $\bar{\bar x}$, $\bar R$, limites de medias y rangos. \\
Kanban & Demanda, lote/contenedor, tiempos de reposicion, defectos & $N_i=D L_i/(y C_i)$ con unidades consistentes. \\
Bodegas & Pallets, pasillos, altura, ciclos de pedido & $d_i^\star=\sqrt{a q_i/(2z_i)}$ y ponderar por frecuencias si hay profundidad unica. \\
Colas & Llegadas por hora, servicio promedio, servidor unico & Calcular $\lambda$, $\mu$, $\rho$, $L_q$, $L$, $W$. \\
Fallas & MTTF, MTTR, disponibilidad & $A=MTTF/(MTTF+MTTR)$ y ajustar tiempo por $1/A$. \\
\bottomrule
\end{longtable}

\begin{materia}{ranking de estudio rapido}
Si queda poco tiempo, el orden recomendado es:
\begin{enumerate}
  \item Guia 11 pregunta 1: muestreo de aceptacion.
  \item Guia 10 pregunta 2: kanban y graficos $\bar X/R$.
  \item Guia 9 problema 1: profundidad de pallets.
  \item Guia 11 pregunta 2: limites $\bar X$ con $\sigma$ conocida.
  \item Guia 8 problemas 1 y 3: M/M/1, CV y disponibilidad.
\end{enumerate}
\end{materia}

\begin{alerta}{como responder para ganar puntos}
Define siempre la variable aleatoria o decision antes de usar formulas. Por ejemplo: ``Sea $X$ el numero de defectuosos en la muestra'', ``sea $N_i$ el numero de kanbans de la operacion $i$'', ``sea $d_i$ la profundidad en pallets del SKU $i$''. Eso evita perder puntos por notacion aunque el calculo este bien.
\end{alerta}
